\section*{Capítulo \ref{ch:b1:electronic-effects} --- Efectos electrónicos e intermedios reactivos}

\begin{solution}{exo:b1:electronic-effects:1}
2-Metilbutano, \ce{CH3-CH(CH3)-CH2-CH3}: los tres \ce{CH3} son primarios;
el \ce{CH2}, secundario, y el CH, terciario. 2-Bromo-2-metilpropano: los
tres \ce{CH3} son primarios, y el carbono central, terciario. De los
cuatro bromuros \ce{C4H9Br} (1-bromobutano, 2-bromobutano,
1-bromo-2-metilpropano, 2-bromo-2-metilpropano), el último da el
\omterm{def:b1:electronic-effects:intermediates}{carbocatión} terciario, el más estable.
\end{solution}

\begin{solution}{exo:b1:electronic-effects:2}
$10^{4.76 - 2.87} = 10^{1.89} = 78$; $10^{4.76 - 0.52} = 10^{4.24} =
\num{1.7e4}$.
\end{solution}

\begin{solution}{exo:b1:electronic-effects:3}
Etanoato: el C=O y el \ce{C-O^-} se intercambian entre los dos oxígenos.
4-Nitrofenóxido: el \omterm{def:b1:lewis-resonance-vsepr:lewis}{par no enlazante} del oxígeno forma un C=O, los
enlaces dobles del anillo se desplazan, el carbono que lleva el \ce{NO2}
forma un enlace C=N y un par N=O pasa al oxígeno: una estructura con C=O
en un extremo del anillo, C=N en el otro y los dos oxígenos del nitro
negativos (una estructura quinoide). Catión alilo:
\ce{CH2=CH-CH2+} $\leftrightarrow$ \ce{+CH2-CH=CH2}.
\end{solution}

\begin{solution}{exo:b1:electronic-effects:4}
Segunda flecha: del enlace O--H del agua a su oxígeno; productos,
\ce{H2O} (procedente del \ce{OH-}, ahora neutro) y \ce{OH-} (el oxígeno
del agua de partida conserva el par, carga $-1$). Para el amoniaco: una
flecha del \omterm{def:b1:lewis-resonance-vsepr:lewis}{par no enlazante} de N a un H del \ce{H3O+} y otra de ese
enlace O--H al oxígeno: N pasa a $+1$, y O, a 0.
\end{solution}

\begin{solution}{exo:b1:electronic-effects:5}
Etanol (15.9) $<$ agua (14.0) $<$ fenol (9.99, deslocalización en el
anillo) $<$ 4-nitrofenol (7.16, el \ce{NO2} toma la carga) $<$ ácido
etanoico (4.76, dos oxígenos equivalentes) $<$ ácido tricloroetanoico
(0.51, tres cloros $-I$).
\end{solution}

\begin{solution}{exo:b1:electronic-effects:6}
Anilina (4.60) $<$ piridina (5.23) $<$ amoniaco (9.25) $<$ metilamina
(10.66). En la anilina, el \omterm{def:b1:lewis-resonance-vsepr:lewis}{par no enlazante} se comparte con el anillo
(estructuras con C=N$^+$ y la carga negativa en \textit{orto} y
\textit{para}); protonar el N costaría esa deslocalización.
\end{solution}

\begin{solution}{exo:b1:electronic-effects:7}
En el 3-nitrofenóxido, las \omterm{def:b1:lewis-resonance-vsepr:resonance}{formas resonantes} sitúan la carga en los
carbonos \textit{orto} y \textit{para} respecto del oxígeno, nunca en el
carbono que lleva el \ce{NO2}: sin ayuda mesómera, de ahí que sea más
débil que el isómero 4. El grupo nitro sigue atrayendo electrones a
través de los enlaces $\sigma$ ($-I$), de ahí que sea más fuerte que el
fenol.
\end{solution}

\begin{solution}{exo:b1:electronic-effects:8}
\ce{CH3+} $<$ \ce{CH3CH2+} $<$ \ce{CH2=CH-CH2+} $<$ \ce{C6H5CH2+}
$\approx$ \ce{(CH3)3C+} $<$ \ce{CH3OCH2+}: los grupos alquilo ceden
electrones; el alilo y el bencilo deslocalizan la carga (dos y cuatro
estructuras); el \omterm{def:b1:lewis-resonance-vsepr:lewis}{par no enlazante} del oxígeno da una estructura con el
octeto completo en todos los átomos. El orden en el centro de la lista
depende del medio.
\end{solution}

\begin{solution}{exo:b1:electronic-effects:9}
\ce{OH-} (\omterm{def:b1:electronic-effects:nucleophile}{nucleófilo}) ataca el carbono de \ce{CH3Br} (\omterm{def:b1:electronic-effects:nucleophile}{electrófilo}), y el
par C--Br se va con el Br. \ce{H2O} cede un par a \ce{H+} (aquí se usan
los nombres de ácido y base). \ce{CN-} ataca el carbono carbonílico del
etanal, y el par $\pi$ del C=O se desplaza al oxígeno. \ce{NH3} cede su
\omterm{def:b1:lewis-resonance-vsepr:lewis}{par no enlazante} al orbital vacío del boro en \ce{BF3}.
\end{solution}

\begin{solution}{exo:b1:electronic-effects:10}
No: $K_a = 10^{-0.51} = 0.31$ es mucho mayor que $C$; el ácido está casi
totalmente disociado. $x^2/(0.010 - x) = 0.31$ da $x = \qty{9.7e-3}{mol/L}$,
$\mathrm{pH} = 2.01$, \omterm{def:b1:predominant-reaction:dissociation}{grado de disociación} del \qty{97}{\percent}.
\end{solution}

\begin{solution}{exo:b1:electronic-effects:11}
Los grupos alquilo empujan electrones ($+I$) hacia el oxígeno del
alcóxido y hacen menos estable su carga, tanto más cuantos más grupos
hay. Además, los alcóxidos más grandes se solvatan peor en el agua, que
estabiliza mejor un oxígeno cargado pequeño que uno oculto entre grupos
alquilo. Ambos efectos elevan el $\mathrm{p}K_a$.
\end{solution}

\begin{solution}{exo:b1:electronic-effects:12}
En una amida, el \omterm{def:b1:lewis-resonance-vsepr:lewis}{par no enlazante} del nitrógeno se comparte con el C=O
(estructura \ce{N+=C-O^-}); cederlo a un protón o a un \omterm{def:b1:electronic-effects:nucleophile}{electrófilo}
destruiría esa estabilización. En \omterm{def:b1:predominant-reaction:strong-acid}{ácido fuerte}, la amida se protona en el
oxígeno, que lleva la carga parcial negativa.
\end{solution}

\begin{solution}{pb:b1:electronic-effects:1}
\textbf{1.} $10^{1.89} = 78$; $10^{1.52} = 33$; $10^{0.84} = 6.9$.
\textbf{2.} No: cada cloro adicional ayuda menos. La carga del carboxilato
ya está bien repartida, y las medidas pierden precisión a medida que el
ácido se acerca a la disociación total.
\textbf{3.} Cada Cl atrae densidad electrónica a través de los enlaces
C--C y C--O y estabiliza el carboxilato ($-I$).
\textbf{4.} El \omterm{def:b1:electronic-effects:inductive}{efecto inductivo} se desvanece en dos o tres enlaces.
\textbf{5.} Los dos ácidos parecen igual de fuertes (0.52 y 0.51). En el
agua, ambos están casi totalmente disociados a las concentraciones
habituales (\cref{ch:b1:predominant-reaction}, \omterm{def:b1:predominant-reaction:levelling}{efecto nivelador}): un
$\mathrm{p}K_a$ cercano a 0 es el límite de lo que el agua puede
distinguir, y los dos valores tienen incertidumbres de unas décimas.
\textbf{6.} Ácido etanoico: forma ácida (\qty{98}{\percent});
cloroetanoico: ambas, con ligero predominio del anión
($10^{3 - 2.87} = 1.3$); los otros tres: anión.
\textbf{7.} \ce{CH3-C(=O)-O^-} $\leftrightarrow$ \ce{CH3-C(-O^-)=O}.
\textbf{8.} Dos enlaces C--O iguales, entre un enlace doble y uno sencillo.
\textbf{9.} La carga está sobre un solo oxígeno, sin enlace $\pi$ con el
que compartirla.
\textbf{10.} $10^{15.9 - 4.76} = 10^{11.1}$, del orden de $10^{11}$.
\textbf{11.} La carga del fenóxido se reparte por tres carbonos del
anillo, menos electronegativos que el oxígeno: una ganancia, menor que en
un carboxilato: $\mathrm{p}K_a$ 9.99, entre 4.76 y 15.9.
\textbf{12.} $K = 10^{6.37 - \mathrm{p}K_a}$: ácido etanoico,
$10^{1.61} = 41$, favorecida (se desprende dióxido de carbono con un
exceso de hidrogenocarbonato); fenol, $10^{-3.62}$, y etanol,
$10^{-9.5}$: no hay reacción.
\textbf{13.} 4-Nitrofenol (7.16) $>$ 2-nitrofenol (7.23) $>$
3-nitrofenol (8.36) $>$ fenol (9.99).
\textbf{14.} La estructura quinoide del
\cref{exo:b1:electronic-effects:3}, con la carga en un oxígeno del grupo
nitro.
\textbf{15.} La carga solo llega a los carbonos \textit{orto} y
\textit{para} respecto del oxígeno; el carbono del nitro está en
\textit{meta}.
\textbf{16.} El efecto $-I$ del \ce{NO2}.
\textbf{17.} En el 2-nitrofenol, el OH forma un \omterm{def:b1:intermolecular-forces-solvents:hbond}{enlace de hidrógeno} con
un oxígeno del grupo nitro vecino, lo que estabiliza la forma ácida y
hace algo más difícil arrancar el protón.
\textbf{18.} Fracciones de anión $1/(1 + 10^{\mathrm{p}K_a - 7.5})$:
4-nitro, \qty{69}{\percent}; 2-nitro, \qty{65}{\percent}; 3-nitro,
\qty{12}{\percent}. El ácido y el anión del 4-nitrofenol absorben de
forma distinta, el anión en el visible: su color aparece en torno a su
$\mathrm{p}K_a$.
\textbf{19.} $x^2/(0.010 - x) = 10^{-0.52} = 0.30$: $x = \qty{9.7e-3}{mol/L}$,
$\mathrm{pH} = 2.01$: casi un \omterm{def:b1:predominant-reaction:strong-acid}{ácido fuerte}.
\textbf{20.} $\frac12(4.76 + 2.00) = 3.38$.
\textbf{21.} \ce{CF3COOH + CH3COO- -> CF3COO- + CH3COOH}, $K = 10^{4.76 -
0.52} = 10^{4.24}$: \omterm{def:b1:extent-q-and-k:quantitative}{cuantitativa}.
\textbf{22.} $K_a(\ce{CF3COOH})/K_a(\ce{CH3COOH}) =$ \textbf{$10^{4.24} =
\num{1.7e4}$}, unas diez mil.
\end{solution}
